\begin{titlingpage}
\begin{minipage}[b]{0.5\textwidth}%
\noindent
\includegraphics[height=100pt,left]{Images/LogoEcolepolytechniquehorizontalbleu.pdf}
\end{minipage}\begin{minipage}[b]{0.5\textwidth}
		\includegraphics[height=100pt,right]{Images/LogoCMLS-Couleur-Version2015-2.pdf}
\end{minipage}

\begin{center}
\begin{otherlanguage}{french}%
\Large{ECOLE DOCTORALE DE L'ECOLE POLYTECHNIQUE}\\[10pt]
 

% Centre de Mathématiques Laurent Schwartz (UMR...)

 
	\Large{\textbf{THÈSE DE DOCTORAT}}

	\large{Discipline: MATHÉMATIQUES}

pour obtenir le grade de

Docteur en Sciences de l'Ecole Polytechnique
 
 présentée par

Timofey ZAKREVSKIY

\end{otherlanguage}%


\noindent\rule{\textwidth}{1pt}
			\LARGE{\bfseries Kinetic models: hydro\-dynamic limits for Fermi-Dirac statistics and solutions of the linearised Boltzmann equation near global Maxwellians.}%
\end{center}
% \\[5pt]

% \noindent\LARGE{\textbf{%
% Modèles cinétiques: sur les limites hydrodynamiques pour la statistique de Fermi-Dirac et (???) sur l'opérateur de scattering pour l'équation de Boltzmann linéarisée.
% }}

\noindent\rule{\textwidth}{1pt}

\vfill

%Caroline's approach
% Date de soutenance: \today
% \\[10pt]
% % Composition de jury: 
% \noindent\begin{tabularx}{\textwidth}{lX >{\bf}c}
% 	&a&b\\
% 	&c&d\\
% 	Composition du jury&a&b\\
% 	&c&d\\
% 	&e&f
% \end{tabularx}

%Isabelle's approach
%\large
{
Soutenue le \today {} devant le jury composée de:

\noindent\begin{tabularx}{\textwidth}{X >{\bfseries}l}
\textbf{François GOLSE}\newline Professeur, Ecole Polytechnique& Directeur de thèse\\
\textbf{John Doe}\newline grade& membre de jury\\
\textbf{John Doe}\newline grade& membre de jury\\
\textbf{John Doe}\newline grade& membre de jury\\
\textbf{John Doe}\newline grade& membre de jury\\
\end{tabularx}
}

\end{titlingpage}